\documentclass[11pt]{article}
\usepackage{amsmath,amssymb,amsthm}

\newtheorem{theorem}{Theorem}

\begin{document}

\begin{theorem}[Napoleon's Theorem]
\label{thm:napoleon}
Let $\triangle ABC$ be any triangle. Construct equilateral triangles all
externally, or all internally, on the three sides of $\triangle ABC$, and let
$L$, $M$, and $N$ be the centroids of the equilateral triangles built on
$BC$, $CA$, and $AB$, respectively. Then $\triangle LMN$ is equilateral.
\end{theorem}

\begin{proof}
Work in vector coordinates in the plane, and let $J$ denote rotation by
$90^\circ$ counterclockwise. There is a sign
$\varepsilon\in\{1,-1\}$ such that
\[
L=\frac{B+C}{2}+\varepsilon\frac{\sqrt{3}}{6}J(C-B),
\]
\[
M=\frac{C+A}{2}+\varepsilon\frac{\sqrt{3}}{6}J(A-C),
\]
\[
N=\frac{A+B}{2}+\varepsilon\frac{\sqrt{3}}{6}J(B-A).
\]
The same sign $\varepsilon$ is used for all three sides: one choice
corresponds to the external construction, and the other to the internal one.

Subtracting gives
\[
M-L=\frac{A-B}{2}+\varepsilon\frac{\sqrt{3}}{6}J(A+B-2C),
\]
\[
N-M=\frac{B-C}{2}+\varepsilon\frac{\sqrt{3}}{6}J(B+C-2A).
\]
Using $J^TJ=I$ and $J^2=-I$, a direct expansion yields
\[
|M-L|^2=|N-M|^2
\]
and
\[
(M-L)\cdot(N-M)=-\frac12|M-L|^2.
\]
Thus the two consecutive side vectors have the same length and make an angle
of $120^\circ$. Equivalently, the interior angle at $M$ is $60^\circ$ and
$LM=MN$.

By the law of cosines,
\[
|L-N|^2
=
|L-M|^2+|M-N|^2-2|L-M||M-N|\cos 60^\circ
=
|L-M|^2.
\]
So $LN=LM=MN$, and therefore $\triangle LMN$ is equilateral.
\end{proof}

\end{document}
